\section{Effective gravitational representation and source-redshift dependence}
\label{sec:representation}

The computed optical response admits a conventional effective-potential
representation. This correspondence is constructive: the geometric and
effective gravitational descriptions encode the same retained light transfer.
The existence of that representation does not identify the physical
preparation from lensing data. Distance-matrix lensing and its associated
degeneracies provide the appropriate comparison
\cite{Fleury2021Arbitrary,Schneider2019Multiplane,SchneiderSluse2014SPT}.

\subsection{Exact correspondence for one source plane}

Use physical angle-to-separation maps in compatible parallel screen frames,
constant across the paraxial beam. Let $T$ be the symmetric time-delay matrix,
$T_0>0$ a scalar reference length, and $M=T/T_0$. The local physical lens
potential $\hat\psi$ and the angular source coordinate
$\boldsymbol\beta=D_{os}^{-1}\boldsymbol X_s$ give
\begin{equation}
 ct(\boldsymbol\theta)
 =\frac12(\boldsymbol\theta-\boldsymbol\beta)^T
 T(\boldsymbol\theta-\boldsymbol\beta)
 -(1+z_l)\hat\psi(D_{ol}\boldsymbol\theta)+f(\boldsymbol\beta).
 \label{eq:representation-original-fermat}
\end{equation}
The arbitrary source-only function $f$ cancels between images of the same
source. Define
\begin{align}
 \boldsymbol\beta_{\rm eff}&=M\boldsymbol\beta,
 \label{eq:representation-effective-source}\\
 \psi_{\rm eff}(\boldsymbol\theta)
 &=\frac{1+z_l}{T_0}\hat\psi(D_{ol}\boldsymbol\theta)
 +\frac12\boldsymbol\theta^T(I-M)\boldsymbol\theta.
 \label{eq:representation-effective-potential}
\end{align}
Expansion of the quadratic terms gives the exact identity
\begin{align}
 \frac{ct(\boldsymbol\theta)}{T_0}
 &=\frac12|\boldsymbol\theta-\boldsymbol\beta_{\rm eff}|^2
   -\psi_{\rm eff}(\boldsymbol\theta)+F_s,\nonumber\\
 F_s&=\frac12\boldsymbol\beta^T(M-M^2)\boldsymbol\beta
       +\frac{f(\boldsymbol\beta)}{T_0}.
 \label{eq:representation-exact-identity}
\end{align}
Since $F_s$ is independent of image position, the stationary images and
relative arrival times are identical in the two representations. The
effective deflection is
\begin{equation}
 \boldsymbol\nabla\psi_{\rm eff}
 =\frac{1+z_l}{T_0}D_{ol}^T
   \hat{\boldsymbol\alpha}(D_{ol}\boldsymbol\theta)
   +(I-M)\boldsymbol\theta.
 \label{eq:representation-effective-deflection}
\end{equation}
This is an application of the established arbitrary-background Fermat
formalism \cite{Fleury2021Arbitrary}; the state-dependent matrices evaluated
here supply its specified cosmological realization.

For the physical SIS potential $\hat\psi(\boldsymbol x)=b|\boldsymbol x|$,
the first term in Eq.~\eqref{eq:representation-effective-potential} is an
elliptical potential when $D_{ol}$ is anisotropic. The quadratic term is a
uniform-convergence and shear contribution. In the nodal quadrupole its
first-order correction is tracefree. Agreement with this effective model
therefore follows from the correspondence and is not separate evidence for
the underlying state. Neither a globally admissible mass distribution nor
agreement with independent dynamical constraints follows from the algebra
alone. Source-position and mass-sheet freedoms must be included when assessing
what an observation can determine \cite{SchneiderSluse2014SPT}.

\subsection{A restricted comparison across two source redshifts}

Consider sources $s$ and $t$ behind the same lens, using the same $D_{ol}$,
physical lens potential, screen frame and small angular patch. For each
source, define background quantities
\begin{align}
 p_s&=\frac{d_{ls}}{d_{os}},&
 T_{0,s}&=(1+z_l)\frac{d_{ol}}{p_s},\nonumber\\
 a_s&=\frac{1+z_l}{T_{0,s}}=\frac{p_s}{d_{ol}},&
 B_s&=I-\frac{T_s}{T_{0,s}}.
 \label{eq:representation-source-scaling}
\end{align}
The equivalent potential has the form
$\psi_{{\rm eff},s}=a_s\hat\psi(D_{ol}\boldsymbol\theta)
+\boldsymbol\theta^TB_s\boldsymbol\theta/2$.
An additional physical quadratic deflector at the main lens,
$\delta\hat\psi(\boldsymbol x)=\boldsymbol x^TA_{\rm lens}\boldsymbol x/2$,
can absorb all these quadratic terms with one fixed symmetric matrix
$A_{\rm lens}$ precisely when
\begin{equation}
 \frac{B_s}{a_s}=D_{ol}^TA_{\rm lens}D_{ol}
 \quad\hbox{is independent of }s.
 \label{eq:representation-physical-quadratic}
\end{equation}
For common nonsingular $D_{ol}$ this is equivalent to source independence
of $B_s/p_s$. To first order about $D_{ol}=d_{ol}I$, it is the familiar
physical-to-reduced scaling $B_s=p_s d_{ol}A_{\rm lens}$.

The tracefree contrast appropriate to the present calculation is
\begin{equation}
 R_{st}=\frac{B_s^{\rm TF}}{p_s}-\frac{B_t^{\rm TF}}{p_t}.
 \label{eq:representation-redshift-contrast}
\end{equation}
With $E_{ab}=-\Gamma_{ab}$ and the common quadrupolar state $y$,
\begin{align}
 B_s^{\rm TF}
   &=[(g_{ol}-g_{ls}+g_{os})\cdot y]H_C+O(|y|^2),\nonumber\\
 R_{st}&=r_{st}H_C+O(|y|^2),\nonumber\\
 r_{st}&=\left[
 \frac{g_{ol}-g_{ls}+g_{os}}{p_s}
 -\frac{g_{ol}-g_{lt}+g_{ot}}{p_t}
 \right]\cdot y.
 \label{eq:representation-state-contrast}
\end{align}
A thin tidal forcing exactly at the main-lens position provides a useful
check: $g_{ol}=g_{ls}=0$ and
$g_{os}\propto S_K(\chi_l)S_K(\chi_s-\chi_l)/S_K(\chi_s)
=S_K(\chi_l)p_s$. Its contrast vanishes, whereas a distributed response
need not obey that scaling.

For the specified nodal tensor
$C_\star=\operatorname{diag}(1,-1,0)/\sqrt2$, the illustrative state with
$y_{\delta_m}=10^{-3}$ and all other initial amplitudes zero gives, at
$z_l=0.3$, $z_s=1$ and $z_t=0.6$,
\begin{equation}
 R_{st}=
 -9.2059866\times10^{-8}
 \begin{pmatrix}1&0\\0&-1\end{pmatrix}
 \quad\hbox{at first order}.
 \label{eq:representation-computed-contrast}
\end{equation}
This dimensionless response cannot be absorbed into the one fixed physical
main-plane quadratic term specified above. It remains compatible with a
distributed conventional gravitational representation of the same light
transfer. The coefficients $B_s$ are model quantities whose inference requires
source, lens and measurement modeling; Eq.~\eqref{eq:representation-computed-contrast}
is not a detected astronomical residual.

In particular, the contrast does not establish general degeneracy breaking.
Nonquadratic lens freedom, source-position transformations, additional lens
planes and uncertain distance factors require their own joint treatment.
Generalized mass-sheet transformations can scale delays for multiple source
planes together \cite{Schneider2019Multiplane}. The result here is narrower:
one shared preparation predicts a definite redshift dependence of its
effective quadratic contribution, and that dependence is constrained when
the physical deflector is required to be shared.
